\def\handoutsay{false}\def\handoutfill{false}
\documentclass[12pt]{article}
\newcommand{\handoutrunning}{Lecture 14: Kronecker Products, vec, and Matrix Derivatives}
\providecommand{\handoutsay}{true}
\usepackage[T1]{fontenc}
\usepackage{lmodern,microtype}
\usepackage[a4paper,margin=15mm,headheight=14pt,headsep=5mm]{geometry}
\usepackage{amsmath,amssymb,amsthm,mathtools}
\usepackage{enumitem,booktabs,needspace}
\usepackage{pifont}
\usepackage[dvipsnames]{xcolor}
\usepackage{fancyhdr}
\usepackage{ifthen}
\usepackage[hidelinks]{hyperref}
\hypersetup{pdfauthor={Chang-Ye Tu}}
\pagestyle{fancy}
\fancyhf{}
\fancyhead[L]{\small Linear Algebra}
\fancyhead[R]{\small\handoutrunning}
\fancyfoot[C]{\small\thepage}
\renewcommand{\headrulewidth}{0.3pt}
\setlength{\parindent}{0pt}
\setlength{\parskip}{4pt plus 1pt minus 1pt}
\setlength{\emergencystretch}{1.5em}
\usepackage{etoolbox}
\newcommand{\handoutdisplayskips}{\setlength{\abovedisplayskip}{7pt plus 2pt minus 2pt}%
  \setlength{\belowdisplayskip}{7pt plus 2pt minus 2pt}}
\apptocmd{\normalsize}{\handoutdisplayskips}{}{\PackageWarning{handout}{could not patch normalsize}}
\handoutdisplayskips
\setlist[enumerate]{label=(\alph*),leftmargin=*,itemsep=2pt,topsep=3pt}
\setlist[itemize]{leftmargin=*,itemsep=2pt,topsep=3pt}
\newcommand{\R}{\mathbb R}
\newcommand{\C}{\mathbb C}
\newcommand{\iu}{\mathrm{i}}
\newcommand{\contradiction}{\mathord{\text{\large\ding{56}}}}
\newcommand{\one}{\mathbf 1}
\DeclareMathOperator{\tr}{tr}
\DeclareMathOperator{\col}{col}
\DeclareMathOperator{\diag}{diag}
\DeclareMathOperator{\rank}{rank}
\DeclareMathOperator\prb{\mathsf{P}}
\DeclareMathOperator\expc{\mathsf{E}}
\DeclareMathOperator\var{var}
\DeclareMathOperator\cov{cov}
\DeclareMathOperator\cor{cor}
\DeclareMathOperator\argmin{argmin}
\DeclareMathOperator\argmax{argmax}
\DeclareMathOperator{\Span}{span}
\theoremstyle{definition}
\newtheorem{theorem}{Theorem}
\newtheorem{proposition}[theorem]{Proposition}
\newtheorem{lemma}[theorem]{Lemma}
\newtheorem{corollary}[theorem]{Corollary}
\newtheorem{definition}[theorem]{Definition}
\newtheorem{example}[theorem]{Example}
\newtheorem{remark}[theorem]{Remark}
\newtheorem{fact}{Fact}
\renewcommand{\thefact}{F\arabic{fact}}
\AddToHook{cmd/subsection/before}{\par\Needspace{9\baselineskip}}
\AddToHook{env/theorem/before}{\par\Needspace{4\baselineskip}}
\AddToHook{env/proposition/before}{\par\Needspace{4\baselineskip}}
\AddToHook{env/lemma/before}{\par\Needspace{4\baselineskip}}
\AddToHook{env/definition/before}{\par\Needspace{4\baselineskip}}
\AddToHook{env/example/before}{\par\Needspace{4\baselineskip}}
\definecolor{saybg}{gray}{0.86}
\NewDocumentEnvironment{say}{+b}
  {\ifthenelse{\equal{\handoutsay}{true}}{\par\Needspace{2\baselineskip}\noindent\colorbox{saybg}{\parbox{\dimexpr\linewidth-2\fboxsep\relax}%
   {\small\sffamily\ding{43}\enspace\textbf{Say.}\enspace#1}}\par}{}}
  {}
\providecommand{\handoutfill}{false}
\newsavebox{\boardbox}
\newlength{\boardht}
\definecolor{boardrule}{gray}{0.55}
\makeatletter
\newcommand{\boardstart}{\par\if@inlabel\leavevmode\par\fi}
\makeatother
\NewDocumentEnvironment{board}{o +b}
  {\boardstart\IfValueT{#1}{\noindent#1\par\nopagebreak}%
   \noindent\begin{lrbox}{\boardbox}\begin{minipage}{\linewidth}#2\end{minipage}\end{lrbox}%
   \setlength{\boardht}{\dimexpr\ht\boardbox+\dp\boardbox\relax}%
   \ifthenelse{\equal{\handoutfill}{true}}%
     {\setlength{\boardht}{1.4\boardht}%
      {\setlength{\fboxsep}{0pt}\setlength{\fboxrule}{0.4pt}\color{boardrule}%
       \framebox[\dimexpr\linewidth-2\fboxrule\relax]{\rule{0pt}{\boardht}}}}%
     {\ifthenelse{\equal{\handoutsay}{true}}%
        {\llap{\color{boardrule}\rule[-\dp\boardbox]{1.5pt}{\boardht}\hspace{4pt}}}{}%
      \usebox{\boardbox}}%
   \par}
  {}
\newcommand{\blank}[1]{\ifthenelse{\equal{\handoutfill}{true}}{\underline{\phantom{#1}}}{#1}}
\newcommand{\lec}[1]{Lecture~#1}
\newcommand{\handouttitle}[2]{%
  \begin{center}
  {\LARGE\sffamily\bfseries #2\par}\vspace{3pt}
  {\sffamily Linear Algebra \textperiodcentered\ \lec{#1}\par}
  \ifthenelse{\equal{\handoutsay}{true}}{\vspace{3pt}{\small\sffamily\colorbox{saybg}{\ding{43}\enspace\textbf{Say.}}\enspace
    teacher's copy: the shaded boxes are not in the students' copy.\par}}{}
  \ifthenelse{\equal{\handoutfill}{true}}{\vspace{3pt}{\small\sffamily Fill-in copy: complete the framed spaces in
    class.\par}}{}
  \end{center}}

\hypersetup{pdftitle={Linear Algebra, Lecture 14: Kronecker Products, vec, and Matrix Derivatives}}
\newcommand{\boardtop}{\vspace*{-\abovedisplayskip}}
\DeclareMathOperator{\row}{row}
\DeclareMathOperator{\vecop}{vec}
\DeclareMathOperator{\adj}{adj}
\begin{document}
\handouttitle{14}{Kronecker Products, vec, and Matrix Derivatives}

\textbf{Last time} (\lec{13}): the singular value decomposition \(A=P\diag(D,0)Q'\), the best approximation of
rank \(k\) (Eckart--Young), principal components through the SVD, the spectral norm and the condition number.
Throughout, \(I\) and \(e_j=\col_j(I)\) are of the size required.

\subsection*{Facts}

From calculus, used without proof (balls, open sets and partial derivatives: Definition~\ref{b16.17}).
\begin{enumerate}[label=\textbf{F\arabic*.},ref=F\arabic*,leftmargin=*,itemsep=0pt,parsep=0pt,topsep=2pt]
\item\label{f-calc} A function on an open set is \(C^1\) (\textbf{continuously differentiable}) if its partial
derivatives exist and are continuous there, \(C^2\) if moreover the partial derivatives of these exist and are
continuous; vector and matrix functions are \(C^1\) (\(C^2\)) if all their entries are. Sums, products and quotients
(nonzero denominator) of \(C^1\) (\(C^2\)) functions are \(C^1\) (\(C^2\)), with
\(\partial(f+g)/\partial x_j=\partial f/\partial x_j+\partial g/\partial x_j\),
\(\partial(fg)/\partial x_j=(\partial f/\partial x_j)g+f\,\partial g/\partial x_j\),
\(\partial(cf)/\partial x_j=c\,\partial f/\partial x_j\) (\(c\) constant); constants and the variables are \(C^2\), so
polynomials in the variables are \(C^2\). A \(C^1\) function is continuous, and a continuous function that is \(\ne0\)
(\(>0\)) at \(c\) is \(\ne0\) (\(>0\)) on some ball around \(c\). \textbf{Chain rule}: if \(g\) is \(C^1\) on an open
\(S\subseteq\R^n\) with values in an open \(T\subseteq\R^m\), and \(f\) is \(C^1\) (\(C^2\)) on \(T\), then \(f\circ g\)
is \(C^1\) on \(S\) (\(C^2\) if \(g\) is too) and \(\mathrm D(f\circ g)(x)=\mathrm Df(g(x))\,\mathrm Dg(x)\). \(\log t\)
(\(t>0\)) is \(C^2\), with derivative \(1/t\).
\item\label{f-sym} If \(f\) is \(C^2\) on an open set, then
\(\partial^2f/\partial x_i\partial x_j=\partial^2f/\partial x_j\partial x_i\), where \(\partial^2f/\partial x_i\partial x_j\)
is the partial derivative with respect to \(x_i\) of \(\partial f/\partial x_j\).
\item\label{f-onevar} (One variable.) \textbf{Fermat's rule}: if \(h\) is differentiable on an open interval and has a
local minimum or maximum at a point \(t_0\) of it, then \(h'(t_0)=0\). \textbf{Taylor's formula}: if \(h\) is \(C^2\)
on an open interval that contains \([0,1]\), then \(h(1)=h(0)+h'(0)+\frac12h''(\tau)\) for some \(0<\tau<1\).
\end{enumerate}

\textbf{From earlier lectures} (restated).
\begin{itemize}[itemsep=0pt,parsep=0pt,topsep=2pt]
\item \lec{1}, Theorem~3(a), Definitions~7, 9, Theorems~8, 10(a),(b): matrix products are associative and
distributive, \(c(AB)=(cA)B=A(cB)\); \((A')'=A\), \((A+B)'=A'+B'\), \((cA)'=cA'\), \((AB)'=B'A'\), \(x'y=y'x\), and a
\(1\times1\) matrix (a number) is its own transpose; \(A\) is \textbf{symmetric} if \(A'=A\); \(\tr A=\sum_ia_{ii}\),
\(\tr(A+B)=\tr A+\tr B\), \(\tr(cA)=c\tr A\), \(\tr(AB)=\tr(BA)\) (\(A\in\R^{m\times n}\),
\(B\in\R^{n\times m}\)). \lec{1}, Definitions~12, 14: \(\diag(d_1,\dots,d_n)\) has the diagonal entries
\(d_1,\dots,d_n\) and all other entries \(0\); the \textbf{summing vector} \(\one_n\in\R^n\) has all entries \(1\).
\item \lec{1}, Theorem~4(a),(b),(d): for \(A=[a_1\ \cdots\ a_n]\in\R^{m\times n}\), \(B=[b_1\ \cdots\ b_p]\in\R^{n\times p}\)
and \(x\in\R^n\): \(\col_j(AB)=Ab_j\); \(Ax=x_1a_1+\dots+x_na_n\), \(Ae_j=a_j\); \((AB)_{ij}=\row_i(A)\,b_j\) and
\(AB=a_1\row_1(B)+\dots+a_n\row_n(B)\) (the \textbf{outer product expansion}). \lec{1}, Proposition~11(a),(b): \(e_i'Ae_j=a_{ij}\); for \(A,B\in\R^{m\times n}\), \(A=B\iff Ax=Bx\) for every
\(x\in\R^n\).
\item \lec{2}, Definition~10, Theorem~15, Corollary~12: \(A\in\R^{n\times n}\) is \textbf{invertible} if
\(AB=BA=I_n\) for some \(B\), the \textbf{inverse} \(A^{-1}\). If \(BA=I\) or \(AB=I\), then \(A\) is invertible and
\(B=A^{-1}\). If \(A\) is symmetric and invertible, \(A^{-1}\) is symmetric.
\item \lec{5}, Definition~1, Theorem~2(b),(c): for a matrix \(A\) partitioned into blocks \(A_{ij}\), \(A'\) has the
\((j,i)\) block \(A_{ij}'\); for a partition of \(B\) conformal for multiplication,
\((AB)_{ik}=A_{i1}B_{1k}+\dots+A_{is}B_{sk}\).
\item \lec{5}, Definition~7, Theorems~12(a), 13(a), 14: \(\det(a)=a\) for \(n=1\); for \(A\in\R^{n\times n}\), \(n\geqslant2\), \(A(i|j)\) is
\(A\) without row \(i\) and column \(j\), \(\det A(i|j)\) a \textbf{minor}, \(c_{ij}=(-1)^{i+j}\det A(i|j)\) a
\textbf{cofactor}, and \(\det A=a_{11}c_{11}+\dots+a_{n1}c_{n1}\). \(A\) is invertible \(\iff\)
\(\det A\ne0\) \(\iff\) \(\rank A=n\); \(\det A=a_{i1}c_{i1}+\dots+a_{in}c_{in}\) for every
\(i\) (expansion along row \(i\)); the \textbf{adjugate} \(\adj A=(c_{ij})'\) has \(A\adj A=(\det A)I_n\), and
\(A^{-1}=\frac1{\det A}\adj A\) if \(\det A\ne0\).
\item \lec{6}, Definition~1: the \textbf{length} \(\lVert x\rVert=\sqrt{x'x}\), the \textbf{distance}
\(\lVert x-y\rVert\); \(\lVert cx\rVert=\lvert c\rvert\,\lVert x\rVert\) (\(c\in\R\)). \lec{13}, Definition~2: the
\textbf{Frobenius norm} \(\lVert A\rVert_F=\bigl(\tr(A'A)\bigr)^{1/2}=\bigl(\sum_i\sum_ja_{ij}^2\bigr)^{1/2}\).
\item \lec{8}, Definition~13, Theorems~14(c), 15(a),(b): a symmetric \(A\) is \textbf{positive definite} (PD) if \(x'Ax>0\) for all \(x\ne0\),
\textbf{positive semidefinite} (PSD) if \(x'Ax\geqslant0\) for all \(x\); a PD matrix is PSD and invertible. If
\(B\in\R^{n\times n}\) is PD and \(K\in\R^{n\times k}\) has \(\rank K=k\), then \(K'BK\) is PD. For
\(A\in\R^{n\times p}\), \(A'A\) is PSD, and PD \(\iff\) \(\rank A=p\).
\item \lec{11}, Definition~1, Theorems~2(b), 3(a): in the linear model \(y=X\beta+\varepsilon\)
(\(X\in\R^{n\times p}\), \(y\in\R^n\)), a \textbf{least-squares estimate} \(\hat\beta\) minimizes
\(S(\beta)=\lVert y-X\beta\rVert^2\); \(X'X\beta=X'y\) are the \textbf{normal equations}. \(X'X\) is invertible \(\iff\)
\(\rank X=p\), and then \(\hat\beta=(X'X)^{-1}X'y\) is their only solution. \(\hat\beta\) is a least-squares estimate
\(\iff\) \(X'X\hat\beta=X'y\). \lec{11}, Example~5: \(X=[\one_4\ (2,0,2,0)'\ (3,3,1,1)']\in\R^{4\times3}\), \(y=(9,3,5,3)'\);
\(X'X=\left(\begin{smallmatrix}4&4&8\\4&8&8\\8&8&20\end{smallmatrix}\right)\) is invertible, so \(\rank X=3\);
\((X'X)^{-1}=\frac14\left(\begin{smallmatrix}6&-1&-2\\-1&1&0\\-2&0&1\end{smallmatrix}\right)\), \(X'y=(20,28,44)'\),
\(\hat\beta=(1,2,1)'\), \(S(\hat\beta)=4\).
\end{itemize}

\subsection*{1. Kronecker products and vec}

\begin{definition}\label{b16.1}
Let \(A=(a_{ij})\in\R^{m\times n}\) and \(B\in\R^{p\times q}\). The \textbf{Kronecker product} of \(A\) and \(B\) is
\[
A\otimes B=\begin{pmatrix}a_{11}B&\cdots&a_{1n}B\\\vdots&&\vdots\\a_{m1}B&\cdots&a_{mn}B\end{pmatrix}\in\R^{mp\times nq}
\]
with \(m\times n\) blocks of size \(p\times q\), the \((i,j)\) block \(a_{ij}B\): the entry in row \((i-1)p+r\) and column
\((j-1)q+s\) is \(a_{ij}b_{rs}\). For \(a\in\R^m\), \(b\in\R^p\): \(a\otimes b=(a_1b',\dots,a_mb')'\in\R^{mp}\). For a
number \(c\), \(c\otimes B=B\otimes c=cB\). Unlike \(AB\), \(A\otimes B\) is defined for any sizes.
\end{definition}

\Needspace{17\baselineskip}
\begin{example}\label{b16.2}
For \(A=(1,\ 2,\ 3)\) and \(B=\begin{pmatrix}6&7\\8&9\end{pmatrix}\), compute \(A\otimes B\), \(B\otimes A\),
\(I_2\otimes\one_3\) and \(\one_3\otimes I_2\).
\end{example}
\textit{Solution.}
\begin{board}[Definition~\ref{b16.1}, block by block: \(A\otimes B=[1B\ \ 2B\ \ 3B]\),
\(B\otimes A=\left(\begin{smallmatrix}6A&7A\\8A&9A\end{smallmatrix}\right)\),
\(I_2\otimes\one_3=\left(\begin{smallmatrix}1\one_3&0\one_3\\0\one_3&1\one_3\end{smallmatrix}\right)\),
\(\one_3\otimes I_2=\left(\begin{smallmatrix}1I_2\\1I_2\\1I_2\end{smallmatrix}\right)\):]
\[
\begin{aligned}
A\otimes B&=\begin{pmatrix}6&7&12&14&18&21\\8&9&16&18&24&27\end{pmatrix}\\[4pt]
B\otimes A&=\begin{pmatrix}6&12&18&7&14&21\\8&16&24&9&18&27\end{pmatrix}
\end{aligned}\qquad
I_2\otimes\one_3=\begin{pmatrix}1&0\\1&0\\1&0\\0&1\\0&1\\0&1\end{pmatrix},\quad
\one_3\otimes I_2=\begin{pmatrix}1&0\\0&1\\1&0\\0&1\\1&0\\0&1\end{pmatrix}
\]
Both \(A\otimes B\) and \(B\otimes A\) contain the twelve numbers \(a_jb_{rs}\), in different orders:
\(A\otimes B\ne B\otimes A\).
\end{board}

\Needspace{10\baselineskip}
\begin{theorem}\label{b16.3}
Let \(A=(a_{ij})\in\R^{m\times n}\) and \(B\in\R^{p\times q}\).
\begin{enumerate}
\item \((A\otimes B)'=A'\otimes B'\).
\item (\textbf{Mixed product}) For \(C=(c_{jk})\in\R^{n\times s}\) and \(D\in\R^{q\times t}\),
\((A\otimes B)(C\otimes D)=(AC)\otimes(BD)\).
\end{enumerate}
\end{theorem}
\begin{proof}
\begin{board}[(a) \(A\otimes B\) has \(m\) block rows and \(n\) block columns; \lec{5}, Theorem~2(b):]
\[
\text{block }(j,i)\text{ of }(A\otimes B)'=(a_{ij}B)'=a_{ij}B'=(A')_{ji}B'=\text{block }(j,i)\text{ of }A'\otimes B'
\]
\end{board}
\begin{board}[(b) \(A\otimes B\) has \(n\) block columns of width \(q\), \(C\otimes D\) has \(n\) block rows of height
\(q\): conformal for multiplication; \lec{5}, Theorem~2(c), and \lec{1}, Theorem~3(a):]
\[
\text{block }(i,k)\text{ of }(A\otimes B)(C\otimes D)=\sum_{j=1}^n(a_{ij}B)(c_{jk}D)=\Bigl(\sum_{j=1}^na_{ij}c_{jk}\Bigr)BD
=(AC)_{ik}BD
\]
This is block \((i,k)\) of \((AC)\otimes(BD)\), for \(i=1,\dots,m\), \(k=1,\dots,s\).
\end{board}
\nopagebreak
\end{proof}

\Needspace{13\baselineskip}
\begin{corollary}\label{b16.4}
\begin{enumerate}
\item \(I_m\otimes I_p=I_{mp}\). For \(B\in\R^{p\times q}\) and \(D=\diag(d_1,\dots,d_m)\), \(D\otimes B\) is block
diagonal with the diagonal blocks \(d_1B,\dots,d_mB\); in particular \(I_m\otimes B\) has \(m\) diagonal blocks \(B\).
\item If \(A\in\R^{m\times m}\) and \(B\in\R^{p\times p}\) are invertible, then \(A\otimes B\) is invertible and
\((A\otimes B)^{-1}=A^{-1}\otimes B^{-1}\).
\item For \(A\in\R^{m\times m}\) and \(B\in\R^{p\times p}\), \(\tr(A\otimes B)=\tr A\cdot\tr B\).
\end{enumerate}
\end{corollary}
\begin{proof}
(a) The \((i,j)\) block of \(D\otimes B\) is \(d_iB\) for \(i=j\) and \(0\) for \(i\ne j\) (Definition~\ref{b16.1};
\lec{1}, Definition~12); \(D=I_m\) gives \(I_m\otimes B\), and then \(B=I_p\) gives \(I_m\otimes I_p=I_{mp}\).
\begin{board}[(b) Theorem~\ref{b16.3}(b); part (a) above:]
\boardtop
\begin{align*}
&(A\otimes B)(A^{-1}\otimes B^{-1})=(AA^{-1})\otimes(BB^{-1})=I_m\otimes I_p=I_{mp}\\
&\implies A\otimes B\text{ is invertible},\ (A\otimes B)^{-1}=A^{-1}\otimes B^{-1}&&\text{(\lec{2}, Theorem~15)}
\end{align*}
\end{board}
\begin{board}[(c) The diagonal blocks of \(A\otimes B\) are \(a_{11}B,\dots,a_{mm}B\); \lec{1}, Theorem~10(a):]
\[
\tr(A\otimes B)=\textstyle\sum_{i=1}^m\tr(a_{ii}B)=\bigl(\sum_{i=1}^ma_{ii}\bigr)\tr B=\tr A\cdot\tr B
\]
\end{board}
\nopagebreak
\end{proof}

\begin{definition}\label{b16.6}
For \(A\in\R^{m\times n}\) with columns \(a_1,\dots,a_n\), \(\vecop A=(a_1',\dots,a_n')'\in\R^{mn}\) stacks the
columns: \(a_{ij}\) is entry \((j-1)m+i\). For instance
\(\vecop\left(\begin{smallmatrix}1&2&3\\4&5&6\end{smallmatrix}\right)=(1,4,2,5,3,6)'\). For \(A,B\in\R^{m\times n}\) and
\(c,d\in\R\): \(\vecop(cA+dB)=c\vecop A+d\vecop B\); \(\vecop A=\vecop B\iff A=B\);
\(\lVert\vecop A\rVert=\lVert A\rVert_F\) (\lec{13}, Definition~2).
\end{definition}

\Needspace{10\baselineskip}
\begin{theorem}\label{b16.7}
\begin{enumerate}
\item For \(a\in\R^m\) and \(b\in\R^n\), \(\vecop(ab')=b\otimes a\).
\item For \(A\in\R^{m\times n}\), \(B\in\R^{n\times p}\) and \(C\in\R^{p\times q}\),
\(\vecop(ABC)=(C'\otimes A)\vecop B\).
\item For \(A,B\in\R^{m\times n}\), \((\vecop A)'\vecop B=\tr(A'B)\).
\end{enumerate}
\end{theorem}
\begin{proof}
\begin{board}[(a) \lec{1}, Theorem~4(a); Definition~\ref{b16.1}:]
\[
\col_j(ab')=b_ja\implies\vecop(ab')=(b_1a',\dots,b_na')'=b\otimes a
\]
\end{board}
\begin{board}[(b) \(b_j=\col_j(B)\), \(u_j=\col_j(I_p)\); \lec{1}, Theorems~4(d), 3(a), 8; then (a) and
Theorem~\ref{b16.3}(b); then the sum over \(j\) (Definition~\ref{b16.6}):]
\boardtop
\begin{align*}
B&=BI_p=\textstyle\sum_{j=1}^pb_ju_j',\qquad ABC=\sum_j(Ab_j)(u_j'C)=\sum_j(Ab_j)(C'u_j)'\\
\vecop\bigl((Ab_j)(C'u_j)'\bigr)&=(C'u_j)\otimes(Ab_j)=(C'\otimes A)(u_j\otimes b_j)=(C'\otimes A)\vecop(b_ju_j')\\
\vecop(ABC)&=(C'\otimes A)\vecop\bigl(\textstyle\sum_jb_ju_j'\bigr)=(C'\otimes A)\vecop B
\end{align*}
\end{board}
\begin{board}[(c) \lec{1}, Theorem~4(d); both vectors list the entries in the same order:]
\[
(A'B)_{jj}=\col_j(A)'\col_j(B)=\textstyle\sum_ia_{ij}b_{ij}\implies\tr(A'B)=\sum_{i,j}a_{ij}b_{ij}=(\vecop A)'\vecop B
\]
\end{board}
\nopagebreak
\end{proof}

\subsection*{2. Derivatives and differentials}

\begin{definition}\label{b16.17}
The \textbf{ball} \(B(c,\delta)\) (\(c\in\R^n\), \(\delta>0\)) is the set of \(x\in\R^n\) with
\(\lVert x-c\rVert<\delta\); \(S\subseteq\R^n\) is \textbf{open} if every \(c\in S\) has a ball \(B(c,\delta)\subseteq
S\). For matrices \(X,C\in\R^{n\times q}\) the distance is \(\lVert X-C\rVert_F=\lVert\vecop X-\vecop C\rVert\)
(Definition~\ref{b16.6}): a set of matrices is open if the set of the vectors \(\vecop X\) is. For a function \(f\) on an
open \(S\) and \(x\in S\), the \textbf{partial derivative} \(\partial f(x)/\partial x_j\) is the derivative at \(t=0\) of
\(t\mapsto f(x+te_j)\), when it exists. For \(f=(f_1,\dots,f_m)':S\to\R^m\) with all these partial derivatives at \(x\),
the \textbf{derivative} (or \textbf{Jacobian matrix}) of \(f\) at \(x\) is
\[
\mathrm Df(x)=\frac{\partial f(x)}{\partial x'}\in\R^{m\times n}\quad\text{with }(i,j)\text{ entry }
\frac{\partial f_i(x)}{\partial x_j}
\]
For \(m=1\), \(\partial f(x)/\partial x=(\mathrm Df(x))'\) is the \textbf{gradient}. A matrix function \(F\), with
\(F(X)\in\R^{m\times p}\), of \(X\) in an open set of \(\R^{n\times q}\) is the vector function
\(\vecop X\mapsto\vecop F(X)\): \(\mathrm DF(X)=\partial\vecop F(X)/\partial(\vecop X)'\in\R^{mp\times nq}\). For a
scalar function \(\varphi(X)\), \(\partial\varphi/\partial X\in\R^{n\times q}\) has the \((i,j)\) entry
\(\partial\varphi/\partial x_{ij}\); so \(\mathrm D\varphi(X)=\bigl(\vecop(\partial\varphi/\partial X)\bigr)'\).
\end{definition}

\begin{definition}\label{b16.20}
Let \(f\) be \(C^1\) (\ref{f-calc}) on an open \(S\subseteq\R^n\), with values in \(\R^m\), and \(x\in S\). The
\textbf{differential} of \(f\) at \(x\) is the linear function \(\mathrm{d}f(x)=\mathrm Df(x)\,\mathrm{d}x\) of an
arbitrary \(\mathrm{d}x=(\mathrm{d}x_1,\dots,\mathrm{d}x_n)'\in\R^n\). For a \(C^1\) matrix function \(F\) of
\(X\in\R^{n\times q}\), \(\mathrm{d}F(X)\) is the matrix of the size of \(F(X)\) with
\(\vecop\mathrm{d}F(X)=\mathrm DF(X)\vecop\mathrm{d}X\) (\(\mathrm{d}X\in\R^{n\times q}\) arbitrary); its \((k,l)\)
entry is \(\sum_{i,j}(\partial f_{kl}/\partial x_{ij})\,\mathrm{d}x_{ij}\). So \(\mathrm{d}f\) has the size of \(f\);
for \(f(x)=x\), \(\mathrm Df=I_n\) and \(\mathrm{d}f=\mathrm{d}x\).
\end{definition}

\Needspace{14\baselineskip}
\begin{theorem}[First identification theorem]\label{b16.21}
\begin{enumerate}
\item Let \(f\) be \(C^1\) on an open \(S\subseteq\R^n\) with values in \(\R^m\), \(x\in S\) and \(A\in\R^{m\times n}\).
Then \(\mathrm{d}f(x)=A\,\mathrm{d}x\) for every \(\mathrm{d}x\) \(\iff\) \(\mathrm Df(x)=A\).
\item Let \(F\) be a \(C^1\) matrix function of \(X\in\R^{n\times q}\) with values in \(\R^{m\times p}\), and
\(A\in\R^{mp\times nq}\). Then \(\vecop\mathrm{d}F(X)=A\vecop\mathrm{d}X\) for every \(\mathrm{d}X\) \(\iff\)
\(\mathrm DF(X)=A\).
\item Let \(\varphi\) be a \(C^1\) scalar function of \(X\in\R^{n\times q}\), and \(C\in\R^{n\times q}\). Then
\(\mathrm{d}\varphi(X)=\tr(C'\mathrm{d}X)\) for every \(\mathrm{d}X\) \(\iff\) \(\partial\varphi/\partial X=C\).
\end{enumerate}
\end{theorem}
\begin{proof}
\begin{board}[(a) Definition~\ref{b16.20}:]
\boardtop
\begin{align*}
&\Leftarrow\colon\ \mathrm{d}f(x)=\mathrm Df(x)\,\mathrm{d}x=A\,\mathrm{d}x\\
&\Rightarrow\colon\ \mathrm Df(x)\,\mathrm{d}x=A\,\mathrm{d}x\text{ for every }\mathrm{d}x\implies\mathrm Df(x)=A
&&\text{(\lec{1}, Proposition~11(b))}
\end{align*}
\end{board}
(b) This is (a) for the vector function \(\vecop X\mapsto\vecop F(X)\) (Definitions~\ref{b16.17}, \ref{b16.20}).
\begin{board}[(c) \(\tr(C'\mathrm{d}X)=(\vecop C)'\vecop\mathrm{d}X\) (Theorem~\ref{b16.7}(c)); (b) with \(m=p=1\),
\(A=(\vecop C)'\); Definition~\ref{b16.17}:]
\[
\mathrm{d}\varphi(X)=(\vecop C)'\vecop\mathrm{d}X\text{ for every }\mathrm{d}X\iff\mathrm D\varphi(X)=(\vecop C)'
\iff\partial\varphi/\partial x_{ij}=c_{ij}\text{ for all }i,j
\]
\end{board}
\nopagebreak
\end{proof}
So, to find a derivative, compute the differential and bring it to one of these forms.

\Needspace{11\baselineskip}
\begin{theorem}\label{b16.22}
Let \(F\), \(G\) be \(C^1\) matrix functions (vector and scalar functions included) of the same variable on an open set,
\(A\) a constant matrix and \(\alpha\) a number.
\begin{enumerate}
\item \(\mathrm{d}A=0\), \(\mathrm{d}(\alpha F)=\alpha\,\mathrm{d}F\), \(\mathrm{d}(F+G)=\mathrm{d}F+\mathrm{d}G\),
\(\mathrm{d}(F')=(\mathrm{d}F)'\), \(\mathrm{d}\vecop F=\vecop\mathrm{d}F\), and \(\mathrm{d}\tr F=\tr\mathrm{d}F\) for
square \(F\).
\item \(\mathrm{d}(FG)=(\mathrm{d}F)G+F\,\mathrm{d}G\) whenever \(FG\) is defined.
\item (\textbf{Chain rule}) Let \(y=g(x)\) be \(C^1\) on an open \(S\subseteq\R^n\) with values in an open
\(T\subseteq\R^m\), and \(z=f(y)\) \(C^1\) on \(T\), with \(\mathrm{d}z=A(y)\,\mathrm{d}y\) and
\(\mathrm{d}y=B(x)\,\mathrm{d}x\). Then the differential of \(x\mapsto f(g(x))\) is
\(\mathrm{d}z=A(g(x))B(x)\,\mathrm{d}x\). The same holds for matrix functions, through \(\vecop\).
\end{enumerate}
\end{theorem}
\begin{proof}
(a) The \((k,l)\) entry of \(\mathrm{d}F\) is \(\sum_{i,j}(\partial f_{kl}/\partial x_{ij})\,\mathrm{d}x_{ij}\)
(Definition~\ref{b16.20}); the partial derivatives of a constant are \(0\), those of \(\alpha f_{kl}\) and
\(f_{kl}+g_{kl}\) are \(\alpha\,\partial f_{kl}/\partial x_{ij}\) and \(\partial f_{kl}/\partial x_{ij}+\partial
g_{kl}/\partial x_{ij}\) (\ref{f-calc}). \(F'\), \(\vecop F\) and \(\tr F\) rearrange or add up entries of \(F\), so their
differentials rearrange or add up the differentials of these entries.
\begin{board}[(b) The product rule of \ref{f-calc}; \((FG)_{kl}=\sum_rf_{kr}g_{rl}\) and (a):]
\boardtop
\begin{align*}
\mathrm{d}(uv)&=\textstyle\sum_{i,j}\bigl((\partial u/\partial x_{ij})v+u\,\partial v/\partial x_{ij}\bigr)\mathrm{d}x_{ij}
=(\mathrm{d}u)v+u\,\mathrm{d}v&&(u,v\text{ scalar})\\
\mathrm{d}(FG)_{kl}&=\textstyle\sum_r\bigl((\mathrm{d}f_{kr})g_{rl}+f_{kr}\,\mathrm{d}g_{rl}\bigr)
=\bigl((\mathrm{d}F)G+F\,\mathrm{d}G\bigr)_{kl}
\end{align*}
\end{board}
(c) By the chain rule (\ref{f-calc}), \(\mathrm D(f\circ g)(x)=\mathrm Df(g(x))\,\mathrm Dg(x)\), and
\(\mathrm Df(y)=A(y)\), \(\mathrm Dg(x)=B(x)\) (Theorem~\ref{b16.21}(a)); so the differential of \(f\circ g\) at \(x\) is
\(A(g(x))B(x)\,\mathrm{d}x\) (Definition~\ref{b16.20}).
\end{proof}

\Needspace{15\baselineskip}
\begin{theorem}\label{b16.24}
Let \(a\in\R^n\) and \(A\in\R^{m\times n}\) be constant.
\begin{enumerate}
\item \(\mathrm{d}(a'x)=a'\,\mathrm{d}x\): \(\mathrm D(a'x)=a'\), and the gradient of \(a'x\) is \(a\).
\item \(\mathrm{d}(Ax)=A\,\mathrm{d}x\): \(\mathrm D(Ax)=A\).
\item For \(A\in\R^{n\times n}\), \(\mathrm{d}(x'Ax)=x'(A+A')\,\mathrm{d}x\): \(\mathrm D(x'Ax)=x'(A+A')\), and the
gradient is \((A+A')x\), which is \(2Ax\) if \(A\) is symmetric.
\item For \(X\in\R^{n\times q}\), \(a\in\R^n\) and \(b\in\R^q\), \(\mathrm{d}(a'Xb)=a'(\mathrm{d}X)b\):
\(\partial(a'Xb)/\partial X=ab'\), and \(\mathrm D(a'Xb)=(b\otimes a)'\).
\end{enumerate}
\end{theorem}
\begin{proof}
\begin{board}[(a), (b) Theorem~\ref{b16.22}(a),(b), then Theorem~\ref{b16.21}(a):]
\boardtop
\begin{align*}
&\text{(b)}\colon\ \mathrm{d}(Ax)=(\mathrm{d}A)x+A\,\mathrm{d}x=A\,\mathrm{d}x\\
&\text{(a)}\colon\ A=a'\text{ in (b)}
\end{align*}
\end{board}
\begin{board}[(c) Theorem~\ref{b16.22}(a),(b); the number \((\mathrm{d}x)'Ax\) is its own transpose (\lec{1},
Theorem~8):]
\boardtop
\begin{align*}
\mathrm{d}(x'Ax)&=(\mathrm{d}x)'Ax+x'A\,\mathrm{d}x=x'A'\,\mathrm{d}x+x'A\,\mathrm{d}x=x'(A+A')\,\mathrm{d}x\\
\mathrm D(x'Ax)&=x'(A+A'),\quad\text{gradient }\bigl(x'(A+A')\bigr)'=(A+A')x&&\text{(Theorem~\ref{b16.21}(a))}
\end{align*}
\end{board}
\begin{board}[(d) Theorem~\ref{b16.22}(a),(b); \lec{1}, Theorem~10(b):]
\boardtop
\begin{align*}
\mathrm{d}(a'Xb)&=a'(\mathrm{d}X)b=\tr\bigl(a'(\mathrm{d}X)b\bigr)=\tr\bigl(ba'\,\mathrm{d}X\bigr)
=\tr\bigl((ab')'\,\mathrm{d}X\bigr)\\
\partial(a'Xb)/\partial X&=ab',\qquad\mathrm D(a'Xb)=\bigl(\vecop(ab')\bigr)'=(b\otimes a)'
&&\text{(Theorems~\ref{b16.21}(c), \ref{b16.7}(a))}
\end{align*}
\end{board}
\nopagebreak
\end{proof}

\Needspace{18\baselineskip}
\begin{theorem}\label{b16.26}
Let \(X\in\R^{n\times n}\) (\(n\geqslant2\)) be a variable and \(A\in\R^{n\times n}\) constant; \(c_{ij}\) are the
cofactors of \(X\) (\lec{5}, Definition~7).
\begin{enumerate}
\item \(\mathrm{d}\tr(AX)=\tr(A\,\mathrm{d}X)\), so \(\partial\tr(AX)/\partial X=A'\).
\item \(\det X\) is a polynomial in the entries of \(X\), \(\partial\det X/\partial x_{ij}=c_{ij}\), and
\(\mathrm{d}\det X=\tr(\adj X\,\mathrm{d}X)\). If \(X\) is invertible, then
\(\mathrm{d}\det X=\det X\,\tr(X^{-1}\mathrm{d}X)\) (\textbf{Jacobi's formula}) and
\(\partial\det X/\partial X=\det X\,(X^{-1})'\).
\item The invertible matrices form an open set, on which \(X\mapsto X^{-1}\) is \(C^2\), and
\(\mathrm{d}X^{-1}=-X^{-1}(\mathrm{d}X)X^{-1}\), \(\mathrm D(X^{-1})=-(X^{-1})'\otimes X^{-1}\).
\item On the open set \(\{X:\det X\ne0\}\), \(\mathrm{d}\log\lvert\det X\rvert=\tr(X^{-1}\mathrm{d}X)\); so
\(\partial\log\det X/\partial X=(X^{-1})'\) where \(\det X>0\).
\end{enumerate}
\end{theorem}
\begin{proof}
\begin{board}[(a) Theorem~\ref{b16.22}(a),(b); \((A')'=A\):]
\[
\mathrm{d}\tr(AX)=\tr\mathrm{d}(AX)=\tr(A\,\mathrm{d}X)=\tr\bigl((A')'\mathrm{d}X\bigr)\implies
\partial\tr(AX)/\partial X=A'\qquad\text{(Theorem~\ref{b16.21}(c))}
\]
\end{board}
\begin{board}[(b) \(\det X\) and its partial derivatives (\lec{5}, Definition~7, Theorem~13(a)):]
\boardtop
\begin{align*}
&\det(x_{11})=x_{11};\quad\det X=\textstyle\sum_kx_{1k}(-1)^{1+k}\det X(1|k)&&\text{(row \(1\))}\\
&\implies\det X\text{ is a polynomial in the entries, so it is \(C^2\)}&&\text{(induction on the size; \ref{f-calc})}\\
&\det X=\textstyle\sum_kx_{ik}c_{ik},\quad c_{ik}=(-1)^{i+k}\det X(i|k)\text{ does not involve row }i&&\text{(row \(i\))}\\
&\implies\det X=c_{ij}x_{ij}+(\text{terms without }x_{ij})\implies\partial\det X/\partial x_{ij}=c_{ij}
\end{align*}
\end{board}
\begin{board}[(b) Definition~\ref{b16.20} (\(1\times1\) case); \((\adj X)_{ji}=c_{ij}\); if \(X\) is invertible,
\(\det X\ne0\) and \(\adj X=\det X\cdot X^{-1}\) (\lec{5}, Theorems~12(a), 14):]
\boardtop
\begin{align*}
\mathrm{d}\det X&=\textstyle\sum_{i,j}c_{ij}\,\mathrm{d}x_{ij}=\sum_j\sum_i(\adj X)_{ji}\,\mathrm{d}x_{ij}
=\sum_j(\adj X\,\mathrm{d}X)_{jj}=\tr(\adj X\,\mathrm{d}X)\\
&=\det X\,\tr(X^{-1}\mathrm{d}X)\quad(X\text{ invertible})\\
\partial\det X/\partial X&=(c_{ij})=(\adj X)'=\det X\,(X^{-1})'\quad(X\text{ invertible})
\end{align*}
\end{board}
\begin{board}[(c) \lec{5}, Theorems~12(a), 14; \ref{f-calc}; Theorem~\ref{b16.22}(a),(b); Theorems~\ref{b16.7}(b),
\ref{b16.21}(b):]
\boardtop
\begin{align*}
&\det C\ne0\implies\det X\ne0\text{ on a ball around }C&&(\det\text{ continuous, (b)})\\
&X^{-1}=\tfrac1{\det X}\adj X\colon\text{ polynomials over a nonzero polynomial}\implies C^2\\
&X^{-1}X=I_n\implies(\mathrm{d}X^{-1})X+X^{-1}\mathrm{d}X=0\implies\mathrm{d}X^{-1}=-X^{-1}(\mathrm{d}X)X^{-1}\\
&\vecop\mathrm{d}X^{-1}=-\bigl((X^{-1})'\otimes X^{-1}\bigr)\vecop\mathrm{d}X\implies\mathrm D(X^{-1})=-(X^{-1})'\otimes X^{-1}
\end{align*}
\end{board}
\begin{board}[(d) \(h(t)=\log t\), \(h'(t)=1/t\) (\ref{f-calc}); Theorem~\ref{b16.22}(c); part (b) above:]
\boardtop
\begin{align*}
\det X>0&\colon\ \mathrm{d}\log\det X=\tfrac1{\det X}\,\mathrm{d}\det X=\tr(X^{-1}\mathrm{d}X)&&\text{(\(\{\det X>0\}\) open: (b), \ref{f-calc})}\\
\det X<0&\colon\ \mathrm{d}\log(-\det X)=\tfrac1{-\det X}(-\mathrm{d}\det X)=\tr(X^{-1}\mathrm{d}X)&&\text{(\(\{\det X<0\}\) open)}\\
\det X>0&\colon\ \partial\log\det X/\partial X=(X^{-1})'&&\text{(Theorem~\ref{b16.21}(c), \(C=(X^{-1})'\))}
\end{align*}
\end{board}
\nopagebreak
\end{proof}

\subsection*{3. The Hessian matrix and extrema}

\Needspace{9\baselineskip}
\begin{definition}\label{b16.31}
Let \(f\) be a \(C^2\) scalar function on an open \(S\subseteq\R^n\). The \textbf{Hessian matrix} of \(f\) at \(x\in S\)
is \(\mathrm Hf(x)=\partial^2f(x)/\partial x\,\partial x'\in\R^{n\times n}\), with \((i,j)\) entry
\(\partial^2f(x)/\partial x_i\partial x_j\); it is symmetric (\ref{f-sym}), and it is the derivative of the gradient,
\(\mathrm Hf=\mathrm D\bigl((\mathrm Df)'\bigr)\). The \textbf{second differential} \(\mathrm{d}^2f(x)\) is the
differential of \(x\mapsto\mathrm{d}f(x)=\mathrm Df(x)\,\mathrm{d}x\) with \(\mathrm{d}x\) held fixed.
\end{definition}
\begin{board}[With \(a_j=\partial f/\partial x_j\); Theorem~\ref{b16.22}(a),(b), \(\mathrm{d}x\) fixed:]
\[
\mathrm{d}f=\textstyle\sum_ja_j\,\mathrm{d}x_j,\qquad
\mathrm{d}^2f=\sum_j(\mathrm{d}a_j)\,\mathrm{d}x_j=\sum_j\sum_i\frac{\partial a_j}{\partial x_i}\,\mathrm{d}x_i\,\mathrm{d}x_j
=(\mathrm{d}x)'\,\mathrm Hf(x)\,\mathrm{d}x
\]
\end{board}
To compute \(\mathrm{d}^2f\), apply Theorem~\ref{b16.22} to \(\mathrm{d}f\), with \(\mathrm{d}(\mathrm{d}x)=0\).

\Needspace{11\baselineskip}
\begin{lemma}\label{b16.19}
\begin{enumerate}
\item For \(A\in\R^{n\times n}\): \(x'Ax=0\) for every \(x\in\R^n\) \(\iff\) \(A'=-A\).
\item For \(A,B\in\R^{n\times n}\): \(x'Ax=x'Bx\) for every \(x\) \(\iff\) \(A+A'=B+B'\). For symmetric \(A\), \(B\),
\(x'Ax=x'Bx\) for every \(x\) only if \(A=B\).
\end{enumerate}
\end{lemma}
\begin{proof}
\begin{board}[(a) A number is its own transpose; \(e_i'Ae_j=a_{ij}\) (\lec{1}, Proposition~11(a)):]
\boardtop
\begin{align*}
&\Leftarrow\colon\ x'Ax=(x'Ax)'=x'A'x=-x'Ax\implies x'Ax=0\\
&\Rightarrow\colon\ x=e_i\colon\ a_{ii}=0;\quad x=e_i+e_j\ (i\ne j)\colon\ 0=a_{ii}+a_{ij}+a_{ji}+a_{jj}=a_{ij}+a_{ji}
\implies A'=-A
\end{align*}
\end{board}
\begin{board}[(b) (a) for \(A-B\):]
\[
x'Ax=x'Bx\ \forall x\iff x'(A-B)x=0\ \forall x\iff(A-B)'=-(A-B)\iff A+A'=B+B'
\]
For symmetric \(A\), \(B\) this reads \(2A=2B\).
\end{board}
\nopagebreak
\end{proof}

\Needspace{9\baselineskip}
\begin{theorem}[Second identification theorem]\label{b16.32}
Let \(f\) be a \(C^2\) scalar function on an open \(S\subseteq\R^n\), \(x\in S\) and \(B\in\R^{n\times n}\). Then
\(\mathrm{d}^2f(x)=(\mathrm{d}x)'B\,\mathrm{d}x\) for every \(\mathrm{d}x\) \(\iff\) \(\mathrm Hf(x)=\frac12(B+B')\).
\end{theorem}
\begin{proof}
\begin{board}[\(\mathrm Hf(x)\) and \(\frac12(B+B')\) are symmetric; \(\mathrm{d}^2f=(\mathrm{d}x)'\mathrm Hf\,\mathrm{d}x\)
(Definition~\ref{b16.31}):]
\boardtop
\begin{align*}
&\Rightarrow\colon\ (\mathrm{d}x)'\mathrm Hf(x)\,\mathrm{d}x=(\mathrm{d}x)'B\,\mathrm{d}x\text{ for every }\mathrm{d}x
\implies2\mathrm Hf(x)=B+B'&&\text{(Lemma~\ref{b16.19}(b))}\\
&\Leftarrow\colon\ (\mathrm{d}x)'\tfrac12(B+B')\,\mathrm{d}x=(\mathrm{d}x)'B\,\mathrm{d}x
&&((\mathrm{d}x)'B'\,\mathrm{d}x=(\mathrm{d}x)'B\,\mathrm{d}x)
\end{align*}
\end{board}
\nopagebreak
\end{proof}
For \(f(x)=x'Ax\) (\(A\in\R^{n\times n}\)): \(\mathrm{d}f=(\mathrm{d}x)'Ax+x'A\,\mathrm{d}x\),
\(\mathrm{d}^2f=2(\mathrm{d}x)'A\,\mathrm{d}x\), so \(\mathrm Hf=A+A'\), which is \(2A\) for symmetric \(A\).

\begin{definition}\label{b16.34}
Let \(f\) be a scalar function on \(S\subseteq\R^n\) and \(c\in S\). \(f\) has a \textbf{local minimum} at \(c\) if
\(f(x)\geqslant f(c)\) for all \(x\in S\) in some ball \(B(c,\delta)\), a \textbf{strict local minimum} if \(f(x)>f(c)\)
there for \(x\ne c\), and a \textbf{global minimum} (\textbf{strict global minimum}) if \(f(x)\geqslant f(c)\) for all
\(x\in S\) (\(f(x)>f(c)\) for \(x\ne c\)); maxima likewise, with the inequalities reversed. If \(S\) is open and \(f\) is
\(C^1\), \(c\) is a \textbf{stationary point} of \(f\) if \(\mathrm Df(c)=0\). \(S\) is \textbf{convex} if
\(tx+(1-t)y\in S\) for all \(x,y\in S\) and \(0\leqslant t\leqslant1\).
\end{definition}

\Needspace{15\baselineskip}
\begin{lemma}\label{b16.35}
Let \(S\subseteq\R^n\) be open, \(f\) a \(C^2\) scalar function on \(S\), \(c\in S\) and \(u\in\R^n\). Let
\(J=\{t\in\R:c+tu\in S\}\) and \(h(t)=f(c+tu)\) for \(t\in J\).
\begin{enumerate}
\item \(J\) is open and contains \(0\); \(h\) is \(C^2\) on \(J\), with \(h'(t)=\mathrm Df(c+tu)u\) and
\(h''(t)=u'\mathrm Hf(c+tu)u\).
\item If \(S\) is convex and \(c+u\in S\), then, for some \(0<\tau<1\),
\(f(c+u)=f(c)+\mathrm Df(c)u+\frac12u'\mathrm Hf(c+\tau u)u\).
\end{enumerate}
\end{lemma}
\begin{proof}
\begin{board}[(a) \(0\in J\) (\(c\in S\)); \lec{6}, Definition~1; the chain rule of \ref{f-calc} (\(t\mapsto c+tu\) has the derivative \(u\)):]
\boardtop
\begin{align*}
&c+t_0u\in S\implies B(c+t_0u,\varepsilon)\subseteq S\ni c+tu\text{ for }\lvert t-t_0\rvert\,\lVert u\rVert<\varepsilon\implies J\text{ open}\\
&h'(t)=\mathrm Df(c+tu)u=\textstyle\sum_ju_j\,\frac{\partial f}{\partial x_j}(c+tu)\\
&h''(t)=\textstyle\sum_j\sum_iu_ju_i\,\frac{\partial^2f}{\partial x_i\partial x_j}(c+tu)=u'\mathrm Hf(c+tu)u,
\text{ continuous in }t&&\text{(Definition~\ref{b16.31})}
\end{align*}
\end{board}
\begin{board}[(b) For \(t\) between \(t_1,t_2\in J\), \(c+tu\) is a convex combination of \(c+t_1u\), \(c+t_2u\):]
\boardtop
\begin{align*}
&J\text{ is an open interval containing \(0\) and \(1\)}&&\text{(\(S\) convex; (a))}\\
&f(c+u)=h(1)=h(0)+h'(0)+\tfrac12h''(\tau)=f(c)+\mathrm Df(c)u+\tfrac12u'\mathrm Hf(c+\tau u)u
&&\text{(\ref{f-onevar}; (a))}
\end{align*}
\end{board}
\nopagebreak
\end{proof}

\Needspace{12\baselineskip}
\begin{theorem}\label{b16.36}
Let \(S\subseteq\R^n\) be open, \(f\) a \(C^2\) scalar function on \(S\), and \(c\in S\).
\begin{enumerate}
\item (\textbf{First-order condition}) If \(f\) has a local minimum or maximum at \(c\), then \(\mathrm Df(c)=0\).
\item (\textbf{Global minimum}) Let \(S\) be convex. If \(\mathrm Df(c)=0\) and \(\mathrm Hf(x)\) is PSD for every
\(x\in S\), then \(f\) has a global minimum at \(c\); if \(\mathrm Hf(x)\) is PD for every \(x\in S\), it is strict.
\item (\textbf{Constraints}) Let \(S\) be convex, \(g\) a \(C^2\) function on \(S\) with values in \(\R^m\),
\(l\in\R^m\), and \(\psi(x)=f(x)-l'g(x)\). If \(\mathrm D\psi(c)=0\), \(g(c)=0\) and \(\mathrm H\psi(x)\) is PSD (PD)
for every \(x\in S\), then \(f(x)\geqslant f(c)\) (\(f(x)>f(c)\) for \(x\ne c\)) for every \(x\in S\) with \(g(x)=0\).
\end{enumerate}
\end{theorem}
\begin{proof}
\begin{board}[(a) Local minimum (for a maximum apply this to \(-f\)): \(f(x)\geqslant f(c)\) on a ball
\(B(c,\delta)\subseteq S\) (\(\delta\) small; \(S\) open); \(h(t)=f(c+te_j)\) (Lemma~\ref{b16.35}, \(u=e_j\)); \(\mathrm Df(c)e_j\) is column \(j\) of \(\mathrm Df(c)\) (\lec{1}, Theorem~4(b)):]
\boardtop
\begin{align*}
&c+te_j\in B(c,\delta)\text{ for }\lvert t\rvert<\delta\implies h\text{ has a local minimum at }t=0\\
&0=h'(0)=\mathrm Df(c)e_j=\partial f(c)/\partial x_j\text{ for every }j\implies\mathrm Df(c)=0
&&\text{(\ref{f-onevar}; Lemma~\ref{b16.35}(a))}
\end{align*}
\end{board}
\begin{board}[(b) \(x\in S\), \(u=x-c\); Lemma~\ref{b16.35}(b) (\(S\) convex, \(c+u=x\in S\)) and \(\mathrm Df(c)=0\):]
\boardtop
\begin{align*}
&f(x)=f(c)+\tfrac12u'\mathrm Hf(c+\tau u)u,\qquad c+\tau u=(1-\tau)c+\tau x\in S\\
&\mathrm Hf(c+\tau u)\text{ PSD}\implies f(x)\geqslant f(c);\qquad\text{PD and }x\ne c\ (u\ne0)\implies f(x)>f(c)
\end{align*}
\end{board}
\begin{board}[(c) \(\psi\) is \(C^2\) on \(S\) (\ref{f-calc}); (b) for \(\psi\):]
\boardtop
\begin{align*}
&\psi(x)\geqslant\psi(c)\text{ for every }x\in S\ (>\text{ for }x\ne c\text{ in the PD case})\\
&g(x)=0\implies f(x)=\psi(x);\qquad g(c)=0\implies\psi(c)=f(c)
\end{align*}
\end{board}
\nopagebreak
\end{proof}
To use (c), solve \(\mathrm D\psi(x)=0\), \(g(x)=0\) for \(x\) and the \textbf{Lagrange multipliers} \(l\)
(\textbf{Lagrange's method}), then check the Hessian of the \textbf{Lagrangian} \(\psi\). For linear constraints
\(g(x)=Rx-r\): \(\mathrm{d}^2(l'Rx-l'r)=\mathrm{d}(l'R\,\mathrm{d}x)=0\), so \(\mathrm H\psi=\mathrm Hf\)
(Theorem~\ref{b16.32}).

\Needspace{13\baselineskip}
\begin{example}\label{b16.37}
Let \(A\in\R^{n\times n}\) be PD, \(b\in\R^n\), and \(f(x)=\frac12x'Ax-b'x\) on \(\R^n\). (a) Minimize \(f\).
(b) Minimize \(f\) subject to \(c'x=0\) (\(c\in\R^n\), \(c\ne0\)).
\end{example}
\textit{Solution.} \(\R^n\) is open and convex; \(f\) is a polynomial, so \(C^2\) (\ref{f-calc}).
\begin{board}[(a) Theorem~\ref{b16.24}(a),(c) with \(A'=A\); Theorem~\ref{b16.32}; \(A\) is invertible (\lec{8},
Theorem~14(c)), \(A^{-1}\) symmetric (\lec{2}, Corollary~12):]
\boardtop
\begin{align*}
&\mathrm{d}f=x'A\,\mathrm{d}x-b'\,\mathrm{d}x=(Ax-b)'\mathrm{d}x,\quad\mathrm{d}^2f=(\mathrm{d}x)'A\,\mathrm{d}x,\quad
\mathrm Hf=A\text{ PD at every }x\\
&\mathrm Df(x)=0\iff Ax=b\iff x=\hat x=A^{-1}b\implies\text{the strict global minimizer}
&&\text{(Theorem~\ref{b16.36}(b))}\\
&(x-\hat x)'A(x-\hat x)=x'Ax-2b'x+b'A^{-1}b=2f(x)+b'A^{-1}b\geqslant0,\ =0\text{ at }x=\hat x\\
&\text{minimum: }f(\hat x)=-\tfrac12b'A^{-1}b
\end{align*}
\end{board}
\begin{board}[(b) \(\psi(x)=f(x)-\lambda c'x\); \lec{5}, Theorem~2(c); \(A^{-1}=(A^{-1})'AA^{-1}\) is PD (\lec{8},
Theorem~15(a), \(K=A^{-1}\), \(\rank A^{-1}=n\) by \lec{5}, Theorem~12(a)):]
\boardtop
\begin{align*}
&\mathrm{d}\psi=(Ax-b-\lambda c)'\mathrm{d}x\colon\ Ax=b+\lambda c,\ c'x=0\iff
\begin{pmatrix}A&c\\c'&0\end{pmatrix}\begin{pmatrix}x\\-\lambda\end{pmatrix}=\begin{pmatrix}b\\0\end{pmatrix}\\
&x=A^{-1}(b+\lambda c),\quad0=c'x=c'A^{-1}b+\lambda c'A^{-1}c,\quad c'A^{-1}c>0\\
&\tilde\lambda=-\frac{c'A^{-1}b}{c'A^{-1}c},\qquad\tilde x=A^{-1}b-\frac{c'A^{-1}b}{c'A^{-1}c}\,A^{-1}c\\
&\mathrm H\psi=\mathrm Hf=A\text{ PD}\implies\tilde x\text{ is the strict minimizer under }c'x=0
&&\text{(Theorem~\ref{b16.36}(c))}
\end{align*}
\end{board}

\Needspace{17\baselineskip}
\subsection*{4. Least squares and the Newton--Raphson method}

\Needspace{14\baselineskip}
\begin{theorem}\label{b16.41}
Let \(X\in\R^{n\times p}\), \(y\in\R^n\), and \(S(\beta)=(y-X\beta)'(y-X\beta)\) for \(\beta\in\R^p\).
\begin{enumerate}
\item \(\mathrm{d}S=-2(y-X\beta)'X\,\mathrm{d}\beta\), so \(\mathrm DS(\beta)=-2(y-X\beta)'X\) and the gradient is
\(2X'X\beta-2X'y\); \(\mathrm HS(\beta)=2X'X\), which is PSD.
\item \(\hat\beta\) minimizes \(S\) \(\iff\) \(X'X\hat\beta=X'y\) (the normal equations).
\item If \(\rank X=p\), then \(\mathrm HS\) is PD, and \(\hat\beta=(X'X)^{-1}X'y\) is the strict global minimizer.
\end{enumerate}
\end{theorem}
\begin{proof}
\begin{board}[(a) \(e=y-X\beta\); Theorems~\ref{b16.22}(a),(b), \ref{b16.24}(b); \((\mathrm{d}e)'e=e'\mathrm{d}e\) (a
number); \(\mathrm{d}\beta\) fixed in \(\mathrm{d}^2S\); \(X'X\) is PSD (\lec{8}, Theorem~15(b)):]
\boardtop
\begin{align*}
&\mathrm{d}e=-X\,\mathrm{d}\beta,\quad\mathrm{d}S=(\mathrm{d}e)'e+e'\,\mathrm{d}e=2e'\,\mathrm{d}e=-2(y-X\beta)'X\,\mathrm{d}\beta\\
&\mathrm DS(\beta)=-2(y-X\beta)'X,\ \text{gradient }-2X'(y-X\beta)=2X'X\beta-2X'y&&\text{(Theorem~\ref{b16.21}(a))}\\
&\mathrm{d}^2S=-2(\mathrm{d}e)'X\,\mathrm{d}\beta=2(\mathrm{d}\beta)'X'X\,\mathrm{d}\beta\implies\mathrm HS=2X'X\text{, PSD}
&&\text{(Theorem~\ref{b16.32})}
\end{align*}
\end{board}
\begin{board}[(b) \(S\) is a polynomial, so \(C^2\) (\ref{f-calc}); \(\R^p\) is open and convex:]
\boardtop
\begin{align*}
&\Rightarrow\colon\ \mathrm DS(\hat\beta)=0\iff X'(y-X\hat\beta)=0&&\text{(Theorem~\ref{b16.36}(a))}\\
&\Leftarrow\colon\ \hat\beta\text{ is a stationary point, }\mathrm HS\text{ PSD everywhere}\implies\text{global minimizer}
&&\text{(Theorem~\ref{b16.36}(b))}
\end{align*}
\end{board}
\begin{board}[(c) \lec{8}, Theorem~15(b); \lec{11}, Theorem~2(b):]
\boardtop
\begin{align*}
&\rank X=p\implies X'X\text{ PD},\ x'(2X'X)x>0\ (x\ne0)\implies\mathrm HS\text{ PD}\\
&X'X\text{ invertible},\ \text{the only solution of the normal equations is }(X'X)^{-1}X'y\\
&\implies\hat\beta=(X'X)^{-1}X'y\text{ is the strict global minimizer}&&\text{(Theorem~\ref{b16.36}(b))}
\end{align*}
\end{board}
\nopagebreak
\end{proof}
Part (b) is \lec{11}, Theorem~3(a), proved there by projection, without calculus.

\Needspace{13\baselineskip}
\begin{example}\label{b16.42}
For \(X\) and \(y\) of \lec{11}, Example~5, find the gradient of \(S\) at \(\beta=0\) and at \(\hat\beta=(1,2,1)'\),
and show that \(\hat\beta\) is the strict global minimizer of \(S\).
\end{example}
\textit{Solution.}
\begin{board}[Theorem~\ref{b16.41}(a),(c): the gradient is \(2X'X\beta-2X'y\); \(X'X\), \(X'y\), \(\rank X=3\) and
\(S(\hat\beta)=4\) from \lec{11}, Example~5:]
\boardtop
\begin{align*}
&X'X=\begin{pmatrix}4&4&8\\4&8&8\\8&8&20\end{pmatrix},\quad X'y=\begin{pmatrix}20\\28\\44\end{pmatrix}\colon\quad
\text{gradient at }\beta=0\text{ is }-2X'y=-2(20,28,44)'\ne0\\
&X'X\hat\beta=(4+8+8,\ 4+16+8,\ 8+16+20)'=(20,28,44)'=X'y\implies\text{gradient at }\hat\beta\text{ is }0\\
&\rank X=3=p\implies\mathrm HS=2X'X\text{ PD}\implies\hat\beta\text{ is the strict global minimizer},\ S(\hat\beta)=4
\end{align*}
\end{board}

\Needspace{21\baselineskip}
\begin{example}[Newton--Raphson]\label{b16.44}
To solve \(n\) equations \(f(x)=0\) in \(n\) unknowns (\(f\) \(C^1\), \(G(x)=\mathrm Df(x)\)), set the first-order
approximation \(f(x_r+\Delta)\approx f(x_r)+G(x_r)\Delta\) (the differential at the current \(x_r\)) to \(0\): the
\textbf{Newton--Raphson} step is
\[
x_{r+1}=x_r-G(x_r)^{-1}f(x_r)\qquad(r=0,1,2,\dots)
\]
provided \(G(x_r)\) is invertible; one stops when the step is negligible. For a stationary point of a \(C^2\) scalar
\(\varphi\): \(f=(\mathrm D\varphi)'\), \(G=\mathrm H\varphi\). Show that one step from any \(x_0\) reaches the
stationary point for \(\varphi(x)=\frac12x'Ax-b'x\) (\(A\) symmetric and invertible) and for the \(S\) of
Theorem~\ref{b16.41} with \(\rank X=p\); take the step from \(\beta_0=0\) in Example~\ref{b16.42}.
\end{example}
\textit{Solution.}
\begin{board}[As in Example~\ref{b16.37}(a), and Theorem~\ref{b16.41}(a),(c); \((X'X)^{-1}\) from \lec{11},
Example~5:]
\boardtop
\begin{align*}
&(\mathrm D\varphi(x))'=Ax-b,\ \mathrm H\varphi=A\implies x_1=x_0-A^{-1}(Ax_0-b)=A^{-1}b\text{, the stationary point}\\
&(\mathrm DS(\beta))'=2X'X\beta-2X'y,\ \mathrm HS=2X'X\text{ invertible}\\
&\implies\beta_1=\beta_0-(2X'X)^{-1}(2X'X\beta_0-2X'y)=(X'X)^{-1}X'y\\
&\beta_0=0\colon\ \beta_1=(X'X)^{-1}X'y
=\frac14\begin{pmatrix}6&-1&-2\\-1&1&0\\-2&0&1\end{pmatrix}\begin{pmatrix}20\\28\\44\end{pmatrix}
=\frac14\begin{pmatrix}120-28-88\\-20+28\\-40+44\end{pmatrix}=\begin{pmatrix}1\\2\\1\end{pmatrix}=\hat\beta
\end{align*}
\end{board}
For other functions several steps are needed.

\bigskip
\textbf{Practice.} The Exercises and Problems of the textbook, Chapter~16: Kronecker products and vec, differentials
and derivatives of traces, determinants and inverses, Hessian matrices and extrema, least squares by differentiation,
the Newton--Raphson method.

\end{document}
